\documentclass[11pt,a4paper]{article}

\usepackage[margin=2.4cm]{geometry}
\usepackage{graphicx}
\usepackage{array}
\usepackage[hidelinks]{hyperref}

\setlength{\parindent}{0pt}
\setlength{\parskip}{0.55em}
\renewcommand{\arraystretch}{1.12}

\title{Assignment 1 -- Poool Game\\Concurrent and Task-Based Implementations}
\author{Alex Santini \\ Francesco Valentini}
\date{PCD a.y. 2025--2026}

\begin{document}

\pagenumbering{gobble}
\begin{titlepage}
    \centering
    \vspace*{2.2cm}

    {\Large PCD -- Concurrent and Distributed Programming\par}
    \vspace{1.2cm}
    {\rule{0.82\textwidth}{0.5pt}\par}
    \vspace{0.9cm}
    {\huge\bfseries Assignment 1\par}
    \vspace{0.35cm}
    {\LARGE Poool Game\par}
    \vspace{0.45cm}
    {\Large Concurrent and Task-Based Implementations\par}
    \vspace{0.9cm}
    {\rule{0.82\textwidth}{0.5pt}\par}

    \vspace{1.6cm}
    {\large\textbf{Authors}\par}
    \vspace{0.25cm}
    {\large Alex Santini\par}
    {\large Francesco Valentini\par}

    \vfill
    {\large Academic Year 2025--2026\par}
\end{titlepage}

\pagenumbering{roman}
\tableofcontents
\newpage
\pagenumbering{arabic}

\section*{Disclaimer}

During the development of the project and the preparation of this report, the authors may have used Large Language Models (LLMs) as auxiliary tools for limited coding support, drafting support, brainstorming, and improving the clarity and structure of the written text. All outputs produced with the assistance of LLMs were critically reviewed, validated, and, when necessary, revised by the authors. All architectural decisions, system design choices, implementation details, and final conclusions were defined and verified by the authors, who take full responsibility for the final content of this report.

\section{Introduction}

Assignment 1 requires three execution strategies for \emph{Poool}: a sequential baseline, a platform-thread implementation, and a task-based implementation built on the Executor framework. The brief also asks for performance measurements and verification with model checking.

The delivered project provides the following execution strategies:

\begin{itemize}
    \item \textbf{Sequential baseline}: single-thread reference execution;
    \item \textbf{Platform-thread version}: Master-Workers architecture;
    \item \textbf{Task-based version}: Executor-based parallel execution.
\end{itemize}

All three launchers reuse the same game rules.

\section{Problem Analysis}

The main difficulty is not the rule set itself, but the combination of shared mutable state, asynchronous activity, and a potentially heavy simulation workload. The board contains many balls that move, collide, and change status during the same tick, so the system must combine concurrency with deterministic state transitions.

\begin{itemize}
    \item \textbf{High computational cost}: collision detection and collision resolution dominate the simulation when the number of balls grows;
    \item \textbf{Shared mutable state}: the simulation must update the state consistently even when several events are active at the same time;
    \item \textbf{Asynchronous inputs}: human and bot actions may arrive while the simulation is running;
    \item \textbf{Deterministic behavior}: concurrency must not change the logical outcome of the match;
    \item \textbf{Responsiveness}: the interface must keep reacting while the simulation progresses.
\end{itemize}

The concurrent part of the assignment therefore becomes a question of identifying which parts of the simulation can be parallelized and which parts must remain serialized. For example, collision processing can be split across independent chunks, while processing pending commands must remain serialized because it requires a consistent global view of the match.

The key challenge is to exploit parallelism without breaking coherence.

\section{Design}

\noindent
The design separates responsibilities to keep the concurrent execution clear. The model owns the rules and the game state, the runtime advances the simulation, and the view reads only stable snapshots.

\begin{itemize}
    \item \textbf{Game logic}: rules and state transitions;
    \item \textbf{Runtime layer}: command handling and tick progression;
    \item \textbf{View layer}: immutable snapshots.
\end{itemize}

The same game logic is reused by the sequential, platform-thread, and task-based versions; only the execution strategy changes.

\subsection{Project architecture}

The project separates what must stay authoritative from what can be scheduled differently. The model keeps the match state and the rules, the runtime layer drives the simulation and coordinates asynchronous actions, and the view consumes only stable snapshots. In this way, the graphical side never depends on the internal timing of the simulation.

\begin{itemize}
    \item the model owns the rules and the authoritative state;
    \item the runtime coordinates commands and tick progression;
    \item the view reads snapshots only, never the live board.
\end{itemize}

The sequential version is the reference architecture. A single execution flow processes commands, advances the physics, updates the logical state, and publishes the resulting snapshot. This version is simple on purpose: it defines the expected semantics and provides the baseline used to evaluate the concurrent versions.

\subsection{Designing the concurrent solution}

The core design choice is to treat the simulation as a mix of local computation and global coordination. The local part, especially collision-related physics work, is where parallelism pays off; the global part remains serialized so that the match evolves through a single coherent sequence of states.

Each tick is therefore divided into phases with different costs and responsibilities. Some phases require a coherent global view of the match and must remain serial: command processing, final state updates, score changes, game-over checks, and snapshot publication. Other phases are suitable for parallel execution because they operate on local portions of the simulation and do not need to modify the shared state at the same time.

\begin{itemize}
    \item \textbf{Parallelizable}: local physics work, especially collision-related computations and other per-chunk updates;
    \item \textbf{Not parallelizable}: command processing, score updates, game-over checks, merge/commit steps, and snapshot publication;
    \item \textbf{Design goal}: keep the coordination points narrow and predictable.
\end{itemize}

This separation is what makes the design effective. The concurrent versions avoid letting multiple threads contend on the same mutable state, and they do not force the whole simulation to wait on fine-grained synchronization. Instead, they keep the coordination points narrow and predictable, while the heavy numerical work runs in parallel on bounded chunks.

\subsection{Platform-thread version}

The platform-thread version adopts a Master-Workers architecture, with the controller acting as the master. The controller keeps ownership of the global progression of the match, while workers handle the parallelizable part of the physics workload. This gives the runtime a clear control point while still allowing the expensive simulation work to run concurrently.

\begin{itemize}
    \item the controller remains responsible for the global tick;
    \item workers handle bounded pieces of the physics workload;
    \item the controller also merges the partial results produced by the workers and applies the final state changes;
    \item the phase boundaries stay explicit, so the protocol remains deterministic.
\end{itemize}

\subsection{Task-based version}

The task-based version keeps a single controller thread for the serial phases of the tick, while the parallelizable physics work is expressed as executor tasks. The controller submits the tasks, waits for completion, merges the partial results, and publishes the final snapshot only after the merge is complete. The difference is mainly in how the concurrent work is represented and scheduled.

\begin{itemize}
    \item the controller keeps the serial phases of the tick;
    \item tasks compute the parallelizable physics work independently;
    \item the controller merges the partial results and publishes the next snapshot only after the merge is complete.
\end{itemize}

\subsection{Shared state and snapshots}

The model creates the logical game snapshot, while the runtime publishes the immutable runtime snapshot consumed by the view. This prevents rendering from seeing half-finished updates and gives the GUI a stable picture of the match at each publication point.

\subsection{Why this design works}

The improvement does not come from changing the game rules, but from reducing coordination overhead around the same rules. As the number of balls grows, the amount of parallelizable physics work grows as well, so the concurrent versions can better exploit the available cores on larger workloads.

\section{Implementation}

\subsection{Main launchers}

\begin{itemize}
    \item \textbf{Sequential launcher}: \texttt{pcd.poool.SequentialPoool};
    \item \textbf{Platform-thread launcher}: \texttt{pcd.poool.ThreadedPoool};
    \item \textbf{Task-based launcher}: \texttt{pcd.poool.TaskBasedPoool}.
\end{itemize}

\subsection{Core modules}

\begin{itemize}
    \item \textbf{\texttt{model/common/math}}: immutable geometric primitives;
    \item \textbf{\texttt{model/physics/common}}: board, balls, holes, and shared physics abstractions;
    \item \textbf{\texttt{model/game}}: game rules and logical snapshots;
    \item \textbf{\texttt{model/physics/sequential}}: baseline stepping logic;
    \item \textbf{\texttt{model/physics/threaded}}: platform-thread engine;
    \item \textbf{\texttt{model/physics/taskbased}}: executor-based engine;
    \item \textbf{\texttt{runtime}}: command queues, command receipts, snapshot publication, bot coordination;
    \item \textbf{\texttt{view}}: Swing-facing rendering classes.
\end{itemize}

\subsection{Important classes}

\begin{itemize}
\item \texttt{GameModel}: owns the game rules and advances the match state;
\item \texttt{Board}: owns the mutable physical state and delegates stepping to a physics engine;
\item \texttt{PhysicsEngine}, \texttt{ThreadedPhysicsEngine}, and \texttt{TaskBasedPhysicsEngine}: implement the simulation strategies;
\item \texttt{ThreadedGameRunner} and \texttt{TaskBasedGameRunner}: coordinate asynchronous commands, controller ticks, and snapshot publication;
\item \texttt{RuntimeGameSnapshot}: packages immutable state for the GUI and asynchronous observers;
\item \texttt{CommandQueueMonitorSupport} and \texttt{SnapshotStoreSupport}: implement the monitor-based coordination used by the concurrent runtime.
\end{itemize}

\section{Verification and Testing}

The repository includes JUnit tests under \texttt{src/test/java/pcd/poool}. They exercise the game model, the physics engines, the runtime helpers, the concurrent runners, the view layer, and the benchmark utilities.

\begin{itemize}
    \item \textbf{JUnit tests}: under \texttt{src/test/java/pcd/poool};
    \item \textbf{Coverage}: model, physics engines, runtime helpers, concurrent runners, view layer, benchmark utilities.
\end{itemize}

\subsection{Petri nets}

The concurrent protocol can be described with an abstract Petri net. The purpose of the model is not to represent the full physics of the game, but to capture one simulation tick as a sequence of controlled phases. In particular, it records command draining, board ownership, worker execution, barrier synchronization, and snapshot publication.

The abstraction is small enough to reason about, yet expressive enough to justify the most relevant safety properties. A tick starts only when the controller owns the board and the command queue has been drained. Worker or task phases then execute on private work units or partial results, and the snapshot is published only after the merge and commit point has been reached. This captures the intended single-writer discipline of the concurrent runtime.

\begin{itemize}
    \item \textbf{What is modeled}: one simulation tick as a sequence of controlled phases;
    \item \textbf{What is not modeled}: full physics, numerical details, GUI rendering;
    \item \textbf{Phases captured}: command draining, board ownership, worker execution, barrier synchronization, snapshot publication;
    \item \textbf{Safety intuition}: the tick starts only when the controller owns the board, and publication happens only after merge and commit.
\end{itemize}

\subsection{JPF model checking}

The repository also includes minimal JPF harnesses in the verification area of the working tree. These harnesses are deliberately smaller than the production runtime: they model a controller, a bounded amount of worker or task activity, a command queue, and the publication step. The goal is to verify the control protocol, not the numerical physics engine.

The JPF models check that commands are not stranded, that publication happens after the tick is complete, and that the reduced protocol can terminate cleanly without deadlock. This is enough to validate the coordination pattern used by both the platform-thread and task-based implementations.

\begin{itemize}
    \item \textbf{Harnesses}: minimal JPF models in the verification area of the working tree;
    \item \textbf{Scope}: controller, bounded worker/task activity, command queue, publication step;
    \item \textbf{Goal}: verify the control protocol, not the numerical physics engine;
    \item \textbf{Properties checked}: no stranded commands, publication after tick completion, clean termination without deadlock.
\end{itemize}

\section{Results and Discussion}

The benchmark snapshot used for this report records results with seed \texttt{42}. The environment was:

\begin{itemize}
    \item \textbf{Seed}: \texttt{42};
    \item \textbf{Environment}: Java 17 on Windows 11, amd64;
    \item \textbf{Visible cores}: 16.
\end{itemize}

The tables below summarize the comparison between the sequential baseline and the concurrent implementations. The first table fixes the worker count and varies the board size; the second fixes the board size and varies the worker count.

\begin{itemize}
    \item \textbf{Speedup table}: compares the sequential median with the concurrent medians across workloads;
    \item \textbf{Worker table}: shows how the two concurrent versions scale as worker count changes;
    \item \textbf{Interpretation}: larger workloads benefit more because the parallelizable physics work outweighs the coordination cost.
\end{itemize}

The speedup snapshot contains the following values:

\begin{table}[htbp]
\centering
\small
\begin{tabular}{rrrrr}
\hline
Balls & Implementation & Seq. median ms & Concurrent median ms & Speedup \\
\hline
1000 & threads  & 138.3033 & 123.2250 & 1.122 \\
1000 & executor & 138.3033 & 122.5104 & 1.129 \\
1500 & threads  & 227.0506 & 153.1364 & 1.483 \\
1500 & executor & 227.0506 & 153.9437 & 1.475 \\
2000 & threads  & 335.7996 & 193.4836 & 1.736 \\
2000 & executor & 335.7996 & 197.4687 & 1.701 \\
2500 & threads  & 450.8448 & 242.6481 & 1.858 \\
2500 & executor & 450.8448 & 237.5317 & 1.898 \\
\hline
\end{tabular}
\caption{Median speedup values from the benchmark snapshot.}
\end{table}

\begin{figure}[htbp]
\centering
\includegraphics[width=0.84\linewidth]{images/speedup-vs-balls.png}
\caption{Speedup as the number of balls increases.}
\end{figure}

The worker-count snapshot contains the following medians for the largest workload in the repository:

\begin{table}[htbp]
\centering
\small
\begin{tabular}{rrrr}
\hline
Implementation & Workers & Median ms & Speedup vs. seq. \\
\hline
threads  & 1  & 1194.7983 & 0.609 \\
threads  & 2  & 760.2068  & 0.957 \\
threads  & 4  & 470.0837  & 1.548 \\
threads  & 8  & 498.5102  & 1.460 \\
threads  & 16 & 427.2735  & 1.703 \\
threads  & 17 & 491.0370  & 1.482 \\
executor & 1  & 1424.2474 & 0.511 \\
executor & 2  & 843.0608  & 0.863 \\
executor & 4  & 582.6597  & 1.249 \\
executor & 8  & 522.3624  & 1.393 \\
executor & 16 & 472.7956  & 1.539 \\
executor & 17 & 442.4130  & 1.645 \\
\hline
\end{tabular}
\caption{Worker-count medians from the benchmark snapshot.}
\end{table}

These numbers also show that the benefit is not monotonic. On the largest workload, the threaded version is slower with 17 workers than with 16, while the executor-based version still improves slightly.

\begin{figure}[htbp]
\centering
\includegraphics[width=0.84\linewidth]{images/speedup-vs-workers.png}
\caption{Speedup as the worker count changes.}
\end{figure}

\section{Conclusions}

The repository contains a sequential implementation, a platform-thread implementation, and a task-based implementation. The shared domain model is reused across all three.

The code uses single-writer ownership, immutable snapshots, explicit synchronization, and a separation between model and runtime.

The benchmark data in the repository shows different speedup values across workloads and worker counts.

\section{How to Run}

The project can be built and tested from the assignment 1 directory. The main launchers are the sequential, platform-thread, and task-based entry points. The report itself can be generated from the \texttt{report} directory.

\begin{itemize}
    \item \textbf{Build and test} from the assignment 1 directory:
\end{itemize}

\begin{verbatim}
mvn -f assignment-1/pom.xml test
\end{verbatim}

\begin{itemize}
    \item \textbf{Main classes}:
\end{itemize}

\begin{verbatim}
pcd.poool.SequentialPoool
pcd.poool.ThreadedPoool
pcd.poool.TaskBasedPoool
\end{verbatim}

\begin{itemize}
    \item \textbf{Report build} from the \texttt{report} directory:
\end{itemize}

\begin{verbatim}
make
\end{verbatim}

\section{References}

No external bibliography items are used. The report is based on the delivered source tree under \texttt{src/main/java/pcd/poool} and \texttt{src/test/java/pcd/poool}.

\end{document}
